%! TEX root = ../main.tex
% =============================================================
% IMMUTABLE — generated automatically, do not edit this block
%   script          : analysis_scratch/all_data_damping_scatter.py
%   plot_type       : damping_all_data_scatter
%   chapter         : 05
%   n_runs_total    : 382
%   conditions      : legacy_nov2025, cond1_h272_high, cond2_h136_high, cond3_h100_high_WRONG, cond4_h100_low
% =============================================================
\begin{figure}[htbp]
  \centering
  \includegraphics[width=\linewidth]{FIGURES/ch05_damping_all_data_scatter.pdf}
  \caption[All-data OUT/IN scatter — cross-condition view]{%
    OUT/IN (FFT) versus $kL$ for all 382 quality-ok full-panel wave runs across the full experimental record. Colour encodes the hardware configuration (probe height $\times$ range mode); marker encodes wind condition (circles = no wind, triangles = full wind); size encodes input amplitude ($0.1$/$0.2$/$0.3$\,V). The blue-shaded band marks the thesis scope (cond4, $1.3$--$1.6$\,Hz) where both wind conditions are present at every frequency. Data outside that band is informative for cross-condition pattern-checking but is not used for the primary results. Extreme outliers (OUT/IN $>2$ or $<0.1$) from low-SNR wind-contamination are clipped from the y-axis for readability; see CLAUDE.md \S16 on the two competing biases in FFT OUT/IN under fullwind.
  }
  \label{fig:ch05_damping_all_data_scatter}
\end{figure}
